\documentclass[11pt]{article}
\usepackage[T1]{fontenc}
\usepackage[utf8]{inputenc}
\usepackage{lmodern,amsmath,amssymb,microtype,xcolor}
\usepackage[margin=.9in]{geometry}
\providecommand{\tightlist}{\setlength{\itemsep}{0pt}\setlength{\parskip}{0pt}}
\usepackage[colorlinks=true,urlcolor=blue!60!black]{hyperref}
\setlength{\parindent}{0pt}
\setlength{\parskip}{.5em}
\setlength{\emergencystretch}{2em}
\author{Wentao Li\\\small Prepared with assistance from OpenAI Codex}
\date{November 22, 2026}
\title{Rainbow Ramsey witnesses and the Boolean prime ideal theorem}
\hypersetup{pdfauthor={Wentao Li},pdftitle={Rainbow Ramsey witnesses and the Boolean prime ideal theorem}}
\begin{document}
\maketitle
Write \(\mathrm{RRT}^r_b(X)\) if every coloring of the \(r\)-element subsets of \(X\) with at most \(b\) sets per color is injective on the \(r\)-element subsets of some infinite \(Y\subseteq X\). A global principle quantifies over all infinite \(X\). Banerjee, Gopaulsingh and Molnár ask whether the Boolean prime ideal theorem (BPI) implies this principle. Venkataramani asks whether the hierarchy in the fiber bound can be strictly increasing at every level in one model, whether systematic arity separations can be constructed, and whether natural restricted rainbow principles can have distinct finiteness classes. \href{https://arxiv.org/abs/2401.02215}{Banerjee, Gopaulsingh and Molnár, 2024, Question 8.2} · \href{https://arxiv.org/abs/2607.22196}{Venkataramani, 2026, Questions 7.15, 7.16, 8.3 and 8.4}

\textbf{Conclusions.} Relative to the consistency of ZFC, the following statements hold.

\begin{enumerate}
\def\labelenumi{\arabic{enumi}.}
\item
  There is a model of pure ZF with BPI in which every global \(\mathrm{RRT}^r_b\), \(r,b\ge2\), fails.
\item
  In another single model of pure ZF, for each \(r,b\ge2\) there is an infinite \(X_{r,b}\) satisfying every \(\mathrm{RRT}^k_s\) for \(2\le k<r\) and finite \(s\), and satisfying \(\mathrm{RRT}^r_s\) exactly for \(1\le s<b\). Thus at fixed arity every increase in the fiber bound gives a strict inclusion of finiteness classes. For fixed bound and any two distinct arities, the lower-arity positive property can hold while the higher one fails.
\item
  Fix an arity \(r\) and allow two inputs with a common color to intersect in any size from a prescribed set \(H\) of integers below \(r\). For colors of multiplicity at most two, these restrictions give finiteness classes ordered exactly by inclusion of the allowed sets \(H\) in a single model of ZF. In particular, at arity two, allowing only disjoint collisions and allowing only collisions at a common vertex give incomparable classes.
\end{enumerate}

\textbf{A class of finite structures.} Fix \(r,b\ge2\). Use a symmetric partial function on \(r\) distinct inputs with an unordered \(b\)-element output. Describe one configuration by disjoint sets \(S,V\) of sizes \(r-1,b\). Declare the function on \(S\cup\{v\}\) to have value \(V\) for each \(v\in V\). Require all other \(r\)-subsets of \(S\cup V\) to be outside its domain. Every function entry must belong to a configuration of this kind. The input and its output recover \(S\) and \(V\) uniquely. The axioms allow configurations to overlap whenever these requirements agree. Repeated-input tuples have no value.

Substructures are closed under the partial function and carry its induced interpretation. The class has free amalgamation over closed substructures. Take the union of the two functions, leaving mixed inputs undefined. Every configuration lies in one side and continues to satisfy its requirements. A new conflicting entry on an old side would have an input entirely in the common substructure and is excluded by closedness and the induced-substructure condition. The class is countable up to isomorphism, is hereditary for closed substructures and admits disjoint joint embeddings. Its countable Fraïssé limit \(M\) is an increasing union of finite closed structures. Every finite set is contained in one. Isomorphisms between finite closed substructures extend to automorphisms, and every finite closed extension embeds over its base. The same construction allows a free linear order, amalgamating the two orders into any common extension.

\textbf{Algebraic closure in the unordered limit.} If \(P\) is closed and finite, every point outside \(P\) has infinite orbit under its pointwise stabilizer. Freely amalgamate arbitrarily many copies of a finite closed extension containing the point over \(P\) and use homogeneity. Every set \(p\) of fewer than \(r\) points is closed. An \(r\)-set with undefined function is closed as well. For an active \(r\)-set \(p=S\cup\{v\}\) its closure is exactly \(S\cup V\) by the restriction on all mixed inputs. The stabilizer of \(p\) acts transitively on the other \(b-1\) petals in \(V\setminus\{v\}\), since every permutation of those petals is an automorphism of that finite closed structure.

\textbf{Witnesses in a permutation model.} Take a disjoint countable sequence of named copies of \(M\) for each pair \((r,b)\). Use their full product automorphism group, which fixes each copy setwise, and the normal filter generated by finite pointwise stabilizers. Let \(X_{r,b}\) be the union of its copies. All constructions in this paragraph take place in the resulting finite-support permutation model of ZFA.

For any coloring with finite support there is a copy disjoint from that support. The equality relation on colors is invariant there under the full automorphism group of that copy, even though the individual colors may move. If two distinct \(k\)-sets \(p,q\) share a color of finite multiplicity, \(q\) has finite orbit under the pointwise stabilizer of \(p\), and conversely. Every vertex of \(q\) therefore has finite orbit under that stabilizer (the union of the orbit of \(q\) is finite). For \(k<r\) closure of \(p\) is \(p\), so \(q\subseteq p\), a contradiction. For \(k=r\) different sets must both be active, and \(q\) must lie in the closure of \(p\). The only active inputs there are the \(b\) sets \(S\cup\{w\}\) for the same \(S,V\). If two of them share a color, invariance under the setwise stabilizer of \(p\) makes the class contain all \(b\) inputs. Thus any coloring with smaller bound is injective on the entire fresh copy.

For failure at bound \(b\), color every active input by its configuration and every inactive input by its own singleton class, with the two types distinguished. These are exactly classes of size \(b\) or one. Any infinite symmetric \(Y\subseteq X_{r,b}\) contains a full configuration. To prove this, take a finite support for \(Y\), enlarge it in its finitely many touched copies to finite closed sets, and choose \(y\in Y\) outside this finite union. In its copy let \(P\) be the closed support and choose a finite closed extension \(T\) containing \(y\). Take \(r-1+b\) copies of \(T\) freely amalgamated over \(P\). On the distinct copies of \(y\), install a single new configuration, using \(r-1\) as the stem and the remaining \(b\) as petals. Each new active input meets more than one copy, and all its newly required inactive mixed inputs were previously undefined. Each copy of \(T\) remains an induced closed substructure. Embed the extension into \(M\) over \(P\). Each chosen vertex is in the orbit of \(y\) under the stabilizer of the support. All belong to \(Y\), so it contains \(b\) same-colored \(r\)-sets. Failure persists for any larger bound.

\textbf{Restricting which intersections can carry the same color.} For \(r\ge2\) and \(H\subseteq\{0,\ldots,r-1\}\), let \(\mathcal F_{r,H}\) consist of sets \(X\) admitting a coloring of \([X]^r\) of multiplicity at most two with no infinite rainbow subset, with the following incidence condition. If distinct \(p,q\) have the same color, then \(|p\cap q|\in H\). These are finiteness classes. They contain finite sets and are closed under taking subsets and bijections. A countably infinite set admits an infinite rainbow subset for every bounded coloring by the ordinary rainbow Ramsey theorem on the natural numbers. The restriction uses only incidence among the inputs and is invariant under relabeling. It gives an answer to Venkataramani's Question 8.4.

When \(H\subseteq K\) we have \(\mathcal F_{r,H}\subseteq\mathcal F_{r,K}\) in ZF. All other inclusions can fail simultaneously. For every \(0\le h<r\) construct a set \(Z_{r,h}\) belonging to \(\mathcal F_{r,H}\) exactly when \(h\in H\). Modify the finite structures above as follows. A configuration consists of two distinct \(r\)-sets \(p,q\) with \(|p\cap q|=h\). Use a symmetric partial function on \(r\) distinct inputs with values of size \(r\). Its value on \(p\) is \(q\) and its value on \(q\) is \(p\). Every entry must be paired in this way, and every other \(r\)-subset of \(p\cup q\) must be inactive. Free amalgamation over a closed induced substructure works by taking the union of the functions. The only new inputs spanning the two sides are left inactive, and no mixed configuration is required. An inactive input is closed and the closure of an active input consists of that input and its mate. Use countably many named copies of the unordered limit for each \((r,h)\), in the same product construction as before. Write \(Z_{r,h}\) for the union of the corresponding copies.

In a fresh copy, suppose two different inputs share a color. Their color class is finite, so by the orbit argument each input lies in the other's closure. An inactive input could therefore share with no distinct input. If \(p\) is active, only \(p\) and its mate are active in its closure, so the other input is precisely its mate. Its intersection with \(p\) has size \(h\). Consequently, if \(h\notin H\), every allowed coloring is injective on a whole fresh copy.

If \(h\in H\), color each pair of mates as one class, and give all other inputs their own singleton classes. Given an infinite supported subset, freely copy the marked finite extension of its support, as in the preceding construction, now \(2r-h\) times. On the distinct marked vertices place the two \(r\)-sets with the specified intersection. Both active inputs span copies because \(r\ge2\). Every other required inactive input on these vertices already was inactive. Homogeneity puts all marked vertices into the given subset, proving failure. For \(H\nsubseteq K\), take \(h\in H\setminus K\) to obtain the promised counterexample. Adding these independent components leaves the arguments for \(X_{r,b}\) intact. Thus at each fixed arity the entire power-set order on allowed intersection sizes is realized. Taking more allowed intersection sizes one at a time also gives chains of arbitrary finite length.

\textbf{Keeping BPI.} Use the freely ordered version of the finite class with \(b=2\). It has the structural Ramsey property by Evans, Hubička and Nešetřil's Theorem 1.3 for free amalgamation classes with partial set-valued functions. Its finite structures are rigid by the order. The Kechris--Pestov--Todorcevic correspondence therefore gives extreme amenability of the automorphism group of its ordered limit. The product over all \(r\ge2\) of the automorphism groups of the ordered copies is again extremely amenable. For instance finite subproducts have the fixed-point property on compact spaces by taking fixed-point sets in succession. The fixed-point sets for the finite-coordinate subproducts have the finite intersection property, and the finite-support subgroup is dense. The product action on its disjoint named components has exactly the product topology. Blass's permutation-model criterion now gives BPI for the model with finite supports. \href{https://arxiv.org/abs/1705.02379}{Evans, Hubička and Nešetřil} · \href{https://arxiv.org/abs/math/0305241}{Kechris, Pestov and Todorcevic} · \href{https://janos.cs.technion.ac.il/RESEARCH/AMS-Book-files/pdfs/11_Blass.pdf}{Blass, Theorems 5.1--5.2}

The negative coloring still works in an ordered copy. In the construction over \(P\) interleave the copies of \(T\) inside each cut of \(P\) consistently with their own orders, so that their distinguished vertices all have the old one-point type over \(P\). Adding the new function values imposes no restrictions on the order. Consequently every infinite symmetric subset contains a collision. The unordered positive results used the symmetry of the petals and are asserted in the unordered model only.

\textbf{Transfer to pure ZF.} The counterexamples and the positive local assertions are boundable. For a fixed base containing the disjoint union of witness sets, their tags and \(\omega\), finite subsets, equivalence relations on those subsets, infinite candidate subsets and injections testing finiteness are all coded within a fixed finite number of power sets. This bound is uniform in \(r,b\) because an unordered finite input is represented as a subset, and a bijection to a natural number is a relation. A coloring with arbitrary codomain enters the property only through its equality relation. The tagged families of witnesses are likewise bounded. For the intersection restrictions, finite inputs, their intersections and the allowed sets of integers are coded at the same bounded rank. The Jech--Sochor embedding transfers the unordered witness statement. For the ordered model use Pincus's transfer theorem for the conjunction of BPI and a boundable (hence injectively boundable) witness statement. This last step preserves BPI in pure ZF. \href{https://doi.org/10.1007/s00605-021-01582-0}{Pincus transfer theorem, statement as Theorem 7}

\textbf{Extent and priority.} Venkataramani established an initial separation with three-point blocks. The construction above independently sets the first active arity and its exact failure threshold. The mate construction further realizes the full inclusion order for restricted intersections. The arity examples supply counterexamples in one direction of the two-parameter problem. It does not decide every reverse arity implication. The ordered version addresses the separate global question about BPI. Provisional significance \textbf{74/100} if correct and new, principally for the separation from a classical weak choice axiom and the uniform family of witnesses. Independent proof and priority review are pending.
\end{document}
